\section{Related Work}
\label{sec:related}

\paragraph{Feed-forward 3D reconstruction.}
DUSt3R~\cite{wang2024dust3r} recasts two-view stereo as direct regression of
aligned pointmaps, removing explicit correspondence search;
MASt3R~\cite{leroy2024mast3r} adds a matching head and metric scale, and
VGGT~\cite{wang2025vggt} generalises the formulation to many views.
Splatt3R~\cite{smart2024splatt3r} builds on MASt3R by freezing the encoder and
decoder and training only a Gaussian head that emits, per pixel, a covariance,
spherical-harmonic colour and opacity (\cref{eq:mu}--\cref{eq:splatt3rloss}).
That freezing decision is not incidental for us: it is the reason the head is
brittle to any change in its input features, which we quantify in
\cref{sec:method:head}. A second consequence matters equally: because
\cref{eq:splatt3rloss} is masked, purely photometric and strictly pairwise, the
training objective is blind to every property that only emerges when hundreds
of keyframes are fused --- which is precisely the regime a SLAM map occupies.

\paragraph{Gaussian-splatting SLAM.}
MonoGS~\cite{matsuki2024monogs} optimises Gaussians and camera poses jointly by
gradient descent from monocular video. Photo-SLAM~\cite{huang2024photoslam}
couples ORB-SLAM3~\cite{campos2021orbslam3} tracking with a
Gaussian mapping thread. SplaTAM~\cite{keetha2024splatam} and
Gaussian-SLAM~\cite{yugay2023gaussianslam} take related routes with RGB-D
input. All of these obtain map primitives by per-scene optimisation, and their
map quality is therefore a function of the optimisation budget they are given
--- a dependency we make explicit when comparing
(\cref{sec:experiments}). Our system differs in that primitives arrive
\emph{predicted}: the optimiser refines an initialisation that is already
approximately correct, rather than creating structure.

\paragraph{Feed-forward priors inside SLAM.}
DROID-SLAM~\cite{teed2021droidslam} established learned dense optimisation for
tracking. MASt3R-SLAM~\cite{murai2025mast3rslam} uses the MASt3R prior for
real-time dense tracking and pointmap fusion, and is the system we build on;
we inherit its tracking front-end unmodified, which is why our trajectory
accuracy matches it and why we are explicit that trajectory accuracy is
\emph{not} a contribution of this work.
VGGT-SLAM~\cite{maggio2025vggtslam} aligns VGGT submaps on the $SL(4)$
manifold; it produces poses and a point cloud but no renderable map, so in our
comparison it serves as a tracking control rather than a rendering baseline.

\paragraph{Parameter-efficient adaptation.}
LoRA~\cite{hu2022lora} is the default tool for adapting a large frozen
backbone. Our experience is a cautionary counterexample: for a pipeline whose
downstream head was \emph{trained against frozen features}, perturbing the
backbone is not a mild intervention, and we measure a $49\%$ collapse. The
useful adaptation surface in such architectures is the head itself.

\paragraph{Evaluation practice.}
Rendering quality in this literature is reported as PSNR/LPIPS~\cite{zhang2018lpips}
and trajectory accuracy as ATE after Sim(3) alignment~\cite{umeyama1991},
typically via evo~\cite{grupp2017evo}. The details underneath those names vary
substantially between papers --- whether scoring views are training or held-out,
whether rendering uses estimated or ground-truth poses, resolution,
undistortion, and post-sequence optimisation budget. \Cref{sec:protocol} shows
that these choices are worth roughly $8.7$\dB{} on unchanged systems, which is
larger than most reported inter-method differences.
